% Prepared on 19 November 2026 as an insertion-ready results subsection.
% Evidence: research/paper1_revision/runs/revision_v1/diagnostics/final_external_v1_20261119T085611Z
% Reproduction: python -m research.paper1_revision.diagnostics.final_external_v1
% Integration also requires updating the abstract, evidence-status section,
% discussion and conclusion, which currently describe an ongoing campaign.

\subsection{Completed independent evaluation and descriptor transfer}
\label{sec:final_independent_descriptors}
The independent campaign completed all 3,067 planned executions on
19 November 2026. Its predictive evaluation comprises 2,000 paired
scenario--deadline observations from 400 new scenarios, with five deadlines
per scenario (5, 10, 30, 60 and 120 seconds). Each of the four batches
contains 100 scenarios: unfiltered draws, altered time structure, modified
topology and decoupled parameters. All 400 reference runs delivered a valid
solution; 389 references were certified optimal. The campaign also completed
217 historical reference upgrades and 450 timing repetitions. The comparisons
below use the predictor archive frozen before independent execution. The
archive's integrity checks passed, and all 2,000 observations were matched to
the saved predictions. The statistical sampling unit is the scenario.

The primary descriptor comparison uses all 22 original raw generator
descriptors or all 22 recomputed MILP descriptors, each standardized using
historical training data. The continuous predictors are the frozen affine
models with nonnegative coefficients; E and ES use their frozen logistic
classifiers. Across all five continuous preferences, the original
unprojected comparison favors the generator descriptors. Realized-descriptor
MSE ranges from 0.40769 to 0.65939, compared with 0.04308 to 0.05618 for the
generator descriptors. The realized affine predictions reach 6.97268 even
though the target is bounded by one.

As a secondary sensitivity analysis, we project the continuous predictions
of every representation onto $[0,1]$. This substantially reduces the realized
models' losses, but generator descriptors retain lower pooled MSE for all
five preferences (Table~\ref{tab:final_independent_descriptors}). Their MSE
is 12.3--41.4\% lower than that of the full realized representation after
projection. These targets are preference-weighted combinations of cost,
unserved energy, elapsed time and direct thermal CO$_2$ intensity. Equal
preference assigns one quarter to each component; each priority assigns
0.6 to the named component and $0.4/3$ to each remaining component.

\begin{table}[htbp]
\centering\small\setlength{\tabcolsep}{4pt}
\begin{tabular}{lrrrrr}
\toprule
Target & Generator & MILP & MILP, B9 limited & Trio & Dimensions\\
\midrule
Equal & 44.35 & 52.75 & 45.73 & 35.25 & 29.23\\
Cost priority & 55.56 & 63.32 & 55.46 & 41.79 & 34.27\\
CO$_2$ priority & 46.01 & 58.65 & 50.42 & 39.62 & 35.50\\
Time priority & 43.07 & 73.52 & 42.00 & 38.31 & 31.14\\
Service priority & 47.38 & 60.54 & 51.86 & 39.80 & 35.65\\
\midrule
Economic success E & 83.38 & 145.27 & 86.28 & 59.93 & 66.27\\
Economic--service success ES & 128.72 & 131.52 & 132.02 & 98.00 & 107.86\\
\bottomrule
\end{tabular}
\caption{Independent-test mean squared error multiplied by $10^3$; lower is
better. Each entry averages 2,000 observations with equal weights. Continuous
predictions are projected onto $[0,1]$; binary entries are Brier losses.
Generator and MILP use 22 descriptors. Trio uses zones, unfixed binary
variables and thermal-zone fraction. Dimensions uses zones, time steps,
lines, variables, unfixed binaries and active constraints. The B9-limited
column is a post-hoc sensitivity that clips the MILP B9 input to its
historical range while retaining every fitted coefficient. It is distinct
from the historical refitted B9 ablation.}
\label{tab:final_independent_descriptors}
\end{table}

For pooled comparisons, 2,000 paired bootstrap resamples draw scenarios
within each batch and retain their five deadlines together. Simultaneous
95\% intervals across the seven generator--MILP target contrasts exclude
zero in favor of generator descriptors for the five projected continuous
targets and for E. For equal preference, the projected MILP-minus-generator
MSE difference is 0.00840, with simultaneous interval [0.00268, 0.01412].
For ES, losses are 0.12872 and 0.13152 for generator and MILP descriptors,
respectively, with difference 0.00280 and interval [$-0.01174$, 0.01735].
The interim ES advantage of the realized representation therefore does not
persist in the complete point estimates, while the final comparison remains
statistically unresolved. Restricting the analysis to the 389 scenarios
with certified references preserves the direction of all seven pooled
generator--MILP differences.

The aggregate ordering varies with the distribution shift. Within the
unfiltered batch, the full realized representation has lower projected MSE
for all five continuous preferences. For equal preference, losses are
0.03127 for MILP and 0.03414 for generator descriptors. Within the topology
batch, the corresponding losses are 0.07198 and 0.04150, and time-priority
losses are 0.16028 and 0.03331. These batch-level statements describe point
estimates. The ordering of the two representations depends on the sampled
distribution shift.

\paragraph{B9 and extrapolation.}
The historical networks are trees, so $B9^{\rm M}=|\mathcal L|/N=1-1/N$;
its training range is [0.75, 0.99412]. In the independent topology batch,
75 scenarios exceed that range, reaching 2.97297. They contribute 375 of
the 2,000 observations; every other independent scenario remains within
the historical B9 range. This is the high-B9 counterpart of the low-B9
extrapolation diagnosed under historical low-criticality family holdout.

To isolate the fitted predictor's sensitivity, we replace B9 by the nearest
endpoint of its historical range only when it lies outside that range,
retain the other inputs and fitted coefficients, and project the resulting
continuous output. Equal-preference MILP MSE falls from 0.05275 to 0.04573,
accounting for 83.6\% of its pooled gap to the generator comparator under
this diagnostic intervention. Time-priority MSE falls from 0.07352 to
0.04200, slightly below the generator value of 0.04307. The intervention
also reduces CO$_2$- and service-priority losses, but residual differences
remain. All five continuous improvements have positive lower bounds in
simultaneous intervals over these five diagnostic contrasts. E Brier loss
falls from 0.14527 to 0.08628, whereas ES slightly worsens from 0.13152 to
0.13202. B9 thus accounts for a substantial, target-dependent part of the
fitted model's transfer error. This decomposition provides predictive
attribution; causal interpretation of network effects remains outside its
scope.

The bounded fractional-logistic regressions and refitted B9 ablations
reported in the historical follow-up remain exploratory development
results. They were outside the pre-execution predictor archive evaluated
here. Independent validation of those fitted alternatives remains
outstanding. Limiting an input at fixed coefficients and refitting a model
after deleting that input address different questions.

\paragraph{Compact predictors.}
The frozen trio of zone count, unfixed binary count and thermal-zone fraction
has 11.1--24.8\% lower projected MSE than the full generator representation
across the five continuous preferences. Its E and ES Brier losses are
0.05993 and 0.09800, reductions of 28.1\% and 23.9\%, respectively.
Simultaneous intervals over its 14 pooled contrasts against the full
generator and full MILP representations favor the trio throughout. The
six-dimensional size representation further reduces continuous MSE by
10.4--18.7\% relative to the trio, with simultaneous intervals over seven
targets favoring it for all five continuous targets. For ES, the trio has
lower loss; for E, their difference remains unresolved. These results
support compact structural predictors as strong comparators, while leaving
room for improvement beyond three indicators.

The pooled summaries, projection and B9 sensitivity were specified after
inspection of interim test outcomes. Their intervals are conditional on
the fitted predictors and adjust within the stated contrast families;
they omit refitting uncertainty and post-selection uncertainty. The
predeclared primary comparison remains the unprojected frozen prediction
comparison by deadline and batch. The completed test uses independent
draws and controlled perturbations from the same generator lineage and
HiGHS implementation. Generalization to other grid datasets, solvers or
hardware remains to be evaluated. Realized feature construction is also
excluded from the archived solve deadlines, so the predictive comparison
alone leaves net online decision benefit unmeasured.
